\section{\BnTeXNativeRegularText{results}}\label{sec:results}
\bnNative{একটি বাংলা শব্দ আগে ব্যবহার করা হয়েছে কি না, তার ওপর এই পদ্ধতির সঠিকতা নির্ভর করে না। নতুন শব্দের জন্যও একই নিয়ম প্রয়োগ করা হয়।}

\subsection*{Sizes and connected forms}
{\small\bnNative{ছোট আকারে: কর্ম, সূর্য, প্রযুক্তি, স্বাধীনতা, স্বাস্থ্য।}\par}
{\normalsize\bnNative{সাধারণ আকারে: রু, রূ, শু, গু, হু, হৃ, ত্রু, ন্তু, স্তু।}\par}
{\large\bnNative{বড় আকারে: ক্ষ, জ্ঞ, হ্ম, ঞ্জ, ঙ্গ, ত্র।}\par}

\begin{figure}[ht]
\centering
\fbox{\parbox{.8\linewidth}{\centering
\vspace{5mm}\bnNative{বাংলা পাঠ ও সাধারণ অনুচ্ছেদের বিন্যাস}\par
\small\bnNative{এই বাক্যটি একটি চিত্রের ভেতরে রাখা হয়েছে।}\vspace{5mm}}}
\caption{\BnTeXNativeText{figure}}\label{fig:example}
\end{figure}

\subsection*{Editing after installation}
\bnNative{পুনর্বিন্যাসযোগ্যতার বিষয়ে নতুন আলোচনা এখানে যোগ করা হলো।}

Edit that sentence or add another explicit Bengali span, then recompile with
pdfLaTeX. Do not regenerate the font installation for a vocabulary change.
The optional memoization package remembers successful glyph runs only within
one compilation; disabling it leaves the same shaping algorithm in use.

Figure~\ref{fig:example} and the bibliography entry~\cite{demo} also exercise
moving or separately included document content.
